\documentclass[11pt]{article}
\usepackage[letterpaper,margin=0.68in]{geometry}
\usepackage{amsmath,amssymb}
\usepackage{enumitem}
\usepackage{xcolor}
\usepackage{fancyhdr}
\usepackage{microtype}
\usepackage{tabularx,array,booktabs}
\usepackage{tikz}
\usepackage[most]{tcolorbox}

\definecolor{olive}{HTML}{4D5430}
\definecolor{gold}{HTML}{9A7B22}
\definecolor{pale}{HTML}{F5F1DC}
\definecolor{softgreen}{HTML}{EDF0E2}
\definecolor{softgray}{HTML}{F5F5F2}
\definecolor{ink}{HTML}{292D20}

\pagestyle{fancy}
\fancyhf{}
\lhead{\textcolor{olive}{MATH\& 107: Math in Society}}
\rhead{\textcolor{gold}{Fall 2026}}
\cfoot{\thepage}
\setlength{\headheight}{14pt}
\setlength{\parindent}{0pt}
\setlength{\parskip}{5pt}
\setlist[enumerate]{leftmargin=*,itemsep=7pt,topsep=5pt}
\setlist[itemize]{leftmargin=1.35em,itemsep=3pt,topsep=3pt}

\newtcolorbox{keybox}[1]{colback=pale,colframe=olive,boxrule=0.8pt,arc=2mm,title=\textbf{#1},fonttitle=\color{olive},left=2mm,right=2mm,top=1.5mm,bottom=1.5mm}
\newtcolorbox{examplebox}[1]{colback=softgreen,colframe=gold,boxrule=0.8pt,arc=2mm,title=\textbf{#1},fonttitle=\color{ink},left=2mm,right=2mm,top=1.5mm,bottom=1.5mm}
\newtcolorbox{refbox}[1]{colback=softgray,colframe=gold,boxrule=0.7pt,arc=2mm,title=\textbf{#1},fonttitle=\color{ink},left=2mm,right=2mm,top=1.5mm,bottom=1.5mm}

\newcommand{\summaryhead}[2]{%
  \clearpage
  {\Large\bfseries\textcolor{olive}{Module #1: #2 — Summary}}\par
  \vspace{2pt}\textcolor{gold}{\rule{\linewidth}{1.2pt}}\par\smallskip
}
\newcommand{\prequizhead}[2]{%
  \clearpage
  {\Large\bfseries\textcolor{olive}{Module #1: #2 — Prequiz}}\par
  \textit{Open-note. Show your mathematical work. A calculation and a clearly stated final answer are enough unless a sentence is specifically requested.}\par
  Name: \hspace{2.6in} Date: \hspace{1.1in}\par
  \vspace{2pt}\textcolor{gold}{\rule{\linewidth}{1.2pt}}\par\smallskip
}
\newcommand{\quizhead}[2]{%
  \clearpage
  {\Large\bfseries\textcolor{olive}{Module #1: #2 — Matching Quiz}}\par
  \textit{These problems use the same skills as the prequiz, with new values and contexts. Show enough work to make your reasoning visible.}\par
  Name: \hspace{2.6in} Date: \hspace{1.1in}\par
  \vspace{2pt}\textcolor{gold}{\rule{\linewidth}{1.2pt}}\par\smallskip
}
\newcommand{\worksmall}{\par\vspace{0.38in}}
\newcommand{\workmed}{\par\vspace{0.62in}}
\newcommand{\worklarge}{\par\vspace{0.92in}}
\newcommand{\workxl}{\par\vspace{1.22in}}
\newcommand{\choice}[1]{\(\square\)~#1\qquad}

\begin{document}
\summaryhead{4}{Intermediate Statistics and Inference}

\textbf{What this module is about.} The normal distribution gives us a way to describe where individual values fall relative to a mean. A z-score measures distance from the mean in standard-deviation units. Statistical inference asks a different question: what can a sample tell us about a population? A confidence interval combines a sample estimate with a margin of error to give a plausible range for a population value.

\begin{keybox}{Normal distributions and the Empirical Rule}
For an approximately normal distribution:
\begin{itemize}
\item about \(68\%\) of values lie within \(1\sigma\) of the mean;
\item about \(95\%\) lie within \(2\sigma\);
\item about \(99.7\%\) lie within \(3\sigma\).
\end{itemize}
The remaining percentage is split equally between the two tails when the distribution is symmetric.
\end{keybox}

\begin{examplebox}{Example 1: Apple weights}
Apple weights are approximately normal with mean \(\mu=180\) g and standard deviation \(\sigma=15\) g.
\begin{itemize}
\item One standard deviation from the mean is \(180\pm15\), or \(165\) to \(195\) g, so about \(68\%\) of apples lie there.
\item Two standard deviations from the mean is \(180\pm30\), or \(150\) to \(210\) g, so about \(95\%\) lie there.
\item About \(5\%\) lie outside that interval, so about \(2.5\%\) weigh more than \(210\) g.
\end{itemize}
\end{examplebox}

\begin{keybox}{z-scores}
For an individual value \(x\),
\[
z=\frac{x-\mu}{\sigma}.
\]
A positive z-score means the value is above the mean; a negative z-score means it is below the mean. To recover a value from a z-score,
\[
x=\mu+z\sigma.
\]
Because z-scores use standard-deviation units, they allow comparisons across different scales.
\end{keybox}

\begin{examplebox}{Example 2: Comparing two standardized tests}
Student A scores \(520\) on a test with mean \(500\) and standard deviation \(100\):
\[
z_A=\frac{520-500}{100}=0.20.
\]
Student B scores \(28\) on a test with mean \(21\) and standard deviation \(5\):
\[
z_B=\frac{28-21}{5}=1.40.
\]
Student B performed farther above the mean relative to the spread of that test.
\end{examplebox}

\begin{keybox}{Samples, parameters, and confidence intervals}
A \textbf{parameter} describes the population. A \textbf{statistic} is computed from a sample and is used to estimate the parameter.

If an estimate is reported as \(E\pm M\), then the corresponding interval is
\[
[E-M,E+M].
\]
A larger sample generally produces a smaller margin of error. In the simplified model used here,
\[
M\propto \frac1{\sqrt n}.
\]
Therefore, multiplying the sample size by \(4\) approximately cuts the margin of error in half.
\end{keybox}

\begin{examplebox}{Example 3: A label claim and a sample interval}
A cereal label claims a population mean of \(230\) calories. A sample gives \(243\pm8\) calories, so the interval is
\[
[243-8,243+8]=[235,251].
\]
Because \(230\) is outside this interval, the label claim is outside the plausible range given by this sample and margin of error.
\end{examplebox}

\prequizhead{4}{Intermediate Statistics and Inference}
\begin{refbox}{Useful formulas}
\[
z=\frac{x-\mu}{\sigma},
\qquad
x=\mu+z\sigma,
\qquad
\text{confidence interval}=\text{estimate}\pm\text{margin of error}.
\]
Empirical Rule: \(68\%\), \(95\%\), \(99.7\%\) within \(1,2,3\) standard deviations.
\end{refbox}

\begin{enumerate}
\item Orange weights are approximately normal with mean \(200\) g and standard deviation \(10\) g.
\begin{enumerate}[label=(\alph*),itemsep=4pt]
\item About what percentage of oranges weigh between \(190\) g and \(210\) g?
\item About what percentage weigh more than \(220\) g?
\item Between what two weights do the middle \(95\%\) of oranges fall?
\end{enumerate}
\worklarge

\item Student A scores \(78\) on a test with mean \(70\) and standard deviation \(4\). Student B scores \(540\) on a different test with mean \(500\) and standard deviation \(50\).
\begin{enumerate}[label=(\alph*),itemsep=4pt]
\item Find each student's z-score.
\item Which student performed better relative to the distribution of scores on that student's test?
\end{enumerate}
\worklarge

\item A consumer group tests bottles of a drink. The sample mean is \(102\) mg of a nutrient per serving, with a reported margin of error of \(\pm5\) mg. The label claims a population mean of \(95\) mg.
\begin{enumerate}[label=(\alph*),itemsep=4pt]
\item Construct the confidence interval.
\item Is the claimed mean of \(95\) mg inside the interval?
\item In this problem, is \(102\) mg a population parameter or a sample statistic?
\end{enumerate}
\workxl

\item A study uses a sample of \(n=150\), and its margin of error is approximately proportional to \(1/\sqrt n\). About what sample size would be needed to cut the margin of error in half? Show how the square-root relationship determines your answer.
\worklarge

\item A distribution has mean \(\mu=50\) and standard deviation \(\sigma=8\). An individual value is \(x=66\).
\begin{enumerate}[label=(\alph*),itemsep=4pt]
\item Find the z-score of \(66\).
\item Is the value within two standard deviations of the mean?
\item Does this z-score answer a question about an individual observation or about a population parameter?
\end{enumerate}
\worklarge
\end{enumerate}

\quizhead{4}{Intermediate Statistics and Inference}
\begin{enumerate}
\item Commute times are approximately normal with mean \(35\) minutes and standard deviation \(8\) minutes. Use the Empirical Rule to estimate the percentage of commutes between \(19\) and \(51\) minutes.
\workmed

\item On one standardized test, a student has \(x=84\), \(\mu=72\), and \(\sigma=6\). Find the z-score and state whether the score is above or below the mean.
\workmed

\item A survey reports \(47\%\) support for a proposal with margin of error \(\pm6\%\). Construct the interval. Is \(55\%\) inside it?
\workmed

\item A sample mean of \(31.2\) cm is computed from \(40\) randomly caught fish. Is \(31.2\) cm a parameter or a statistic? What population quantity is it intended to estimate?
\workmed

\item If a sample size is multiplied by \(9\), by what factor does the simplified \(1/\sqrt n\) model predict the margin of error will change?
\workmed
\end{enumerate}

\end{document}
